% Claim statement file for row claim_3_17 -- "MarginalizationsCommute".
% Auto-generated stub by scaffold/solve_chapter.py at row start. The
% manager agent (or a worker it dispatches) fills the body below with the
% claim's statement(s) from the lecture notes -- statement only, no proof
% (proofs live in the sibling `claim_3_17_proof_*.tex` file).
%
% This file is a sibling of `main.tex` inside `<subsection>/tex/`. The
% `subfiles` package lets it render both standalone and as part of
% `main.tex`. Note: `main.tex` does NOT \subfile this statement file -- the
% sibling `_proof_` file restates the statement above the proof, so
% including both in `main.tex` would be redundant. This file exists for
% standalone reading / grep convenience.

\documentclass[main]{subfiles}

\begin{document}
% ref: claim_3_17 (statement)
% title: MarginalizationsCommute
\def\rowref{claim\_3\_17}
\phantomsection\label{claim_3_17}

\begin{Lem}[Marginalizations commute]\label{marginalizations-commute}
    Let $G = (J, V, E, L)$ be a CDMG (def \ref{def-cdmg}); throughout we use the notation of def \ref{not-cdmg} and recall the typing and axioms of def \ref{def-cdmg}, namely $J \cap V = \emptyset$, $E \ins (J \cup V) \times V$, $L \ins V \times V$, $L$ irreflexive ($(v_1, v_2) \in L \Rightarrow v_1 \neq v_2$), and $L$ symmetric ($(v_1, v_2) \in L \Leftrightarrow (v_2, v_1) \in L$). We rely throughout on the marginalization operator $G^{\sm W}$ of def \ref{def:G_marginalization}. Let
    \[
        W_1 \;\ins\; V \quad\text{and}\quad W_2 \;\ins\; V
    \]
    be two subsets of output nodes of $G$ (both subsets are constrained to lie on the output side $V$, matching the precondition ``$W \ins V$'' of def \ref{def:G_marginalization} for the inner marginalizations $G^{\sm W_1}$ and $G^{\sm W_2}$), subject to the disjointness side condition
    \[
        W_1 \cap W_2 \;=\; \emptyset.
    \]
    All three quantifiers ``for every CDMG $G$'', ``for every $W_1 \ins V$'', and ``for every $W_2 \ins V$ with $W_1 \cap W_2 = \emptyset$'' are read \emph{universally}; in particular the (vacuously disjoint) degenerate cases $W_1 = W_2 = \emptyset$, $W_1 = \emptyset$ alone, and $W_2 = \emptyset$ alone are all admitted (def \ref{def:G_marginalization} explicitly admits the $W = \emptyset$ input, in which case $G^{\sm \emptyset} = G$ verbatim by items~i.--iv.\ of that def, and the displayed triple equality below collapses accordingly).

    \emph{Well-typedness of the iterated and joint marginalizations.} Each of the three CDMGs appearing in the displayed triple equality below is well-typed per def \ref{def:G_marginalization}'s precondition ``$W \ins V$ of the input CDMG'', evaluated against the input CDMG's output-node set (item~ii.\ of that def):
    \begin{itemize}
        \item \emph{LHS, $\lp G^{\sm W_1} \rp^{\sm W_2}$.} The inner $G^{\sm W_1}$ is defined per def \ref{def:G_marginalization} from $W_1 \ins V$. By items~i.\ and~ii.\ of that def, the input-node and output-node sets of $G^{\sm W_1}$ are
              \[ J^{\sm W_1} \;=\; J \qquad\text{and}\qquad V^{\sm W_1} \;=\; V \sm W_1. \]
              For the outer marginalization $(\,\cdot\,)^{\sm W_2}$ applied to $G^{\sm W_1}$ to be defined per def \ref{def:G_marginalization}, we need $W_2$ to be a subset of the output-node set $V^{\sm W_1} = V \sm W_1$ of $G^{\sm W_1}$. This holds because $W_2 \ins V$ together with $W_1 \cap W_2 = \emptyset$ yields $W_2 \ins V \sm W_1 = V^{\sm W_1}$. \emph{This is precisely the load-bearing role of the disjointness hypothesis $W_1 \cap W_2 = \emptyset$:} without it, $W_2$ could overlap with $W_1$ -- a set of nodes that has already been removed from the output side by the inner marginalization $G^{\sm W_1}$ -- and the outer marginalization on $G^{\sm W_1}$ would be ill-typed against def \ref{def:G_marginalization}'s precondition.

        \item \emph{Mirror, $\lp G^{\sm W_2} \rp^{\sm W_1}$.} Symmetrically: the inner $G^{\sm W_2}$ is defined per def \ref{def:G_marginalization} from $W_2 \ins V$, with input-node and output-node sets
              \[ J^{\sm W_2} \;=\; J \qquad\text{and}\qquad V^{\sm W_2} \;=\; V \sm W_2 \]
              (items~i.\ and~ii.\ of def \ref{def:G_marginalization}). The outer marginalization $(\,\cdot\,)^{\sm W_1}$ applied to $G^{\sm W_2}$ is defined because $W_1 \ins V$ together with $W_1 \cap W_2 = \emptyset$ yields $W_1 \ins V \sm W_2 = V^{\sm W_2}$.

        \item \emph{Joint, $G^{\sm (W_1 \cup W_2)}$.} The single marginalization $G^{\sm (W_1 \cup W_2)}$ is defined per def \ref{def:G_marginalization} from $W_1 \cup W_2 \ins V$, which is immediate from $W_1 \ins V$ and $W_2 \ins V$. (Disjointness is \emph{not} needed to verify the subset-of-$V$ precondition on the joint side, but is built into the lemma's hypothesis so that all three CDMGs in the triple equality are simultaneously well-typed.)
    \end{itemize}

    \emph{Carrier identification across the three CDMGs.} By items~i.\ and~ii.\ of def \ref{def:G_marginalization} applied iteratively (for the two iterated marginalizations) and once (for the joint marginalization), the input-node and output-node sets of the three CDMGs in the displayed triple equality below coincide:
    \begin{align*}
        J_{\lp G^{\sm W_1} \rp^{\sm W_2}} \;&=\; J_{\lp G^{\sm W_2} \rp^{\sm W_1}} \;=\; J_{G^{\sm (W_1 \cup W_2)}} \;=\; J, \\
        V_{\lp G^{\sm W_1} \rp^{\sm W_2}} \;&=\; V_{\lp G^{\sm W_2} \rp^{\sm W_1}} \;=\; V_{G^{\sm (W_1 \cup W_2)}} \;=\; V \sm (W_1 \cup W_2).
    \end{align*}
    The $J$-component identification follows because every application of def \ref{def:G_marginalization} fixes $J$ (item~i.: $J^{\sm W} = J$); the $V$-component identification uses item~ii.\ of def \ref{def:G_marginalization} together with the elementary set identities $(V \sm W_1) \sm W_2 = V \sm (W_1 \cup W_2)$ and $(V \sm W_2) \sm W_1 = V \sm (W_1 \cup W_2)$ on the iterated marginalizations, and direct substitution of $W_1 \cup W_2$ into item~ii.\ on the joint marginalization. Consequently, the asserted triple equality of CDMGs is well-typed at the level of input/output partitions, and the substantive content of the lemma -- deferred to the proof -- is the equality of the directed-edge and bidirected-edge sets.

    \emph{The conclusion.} The three CDMGs $\lp G^{\sm W_1} \rp^{\sm W_2}$, $\lp G^{\sm W_2} \rp^{\sm W_1}$, and $G^{\sm (W_1 \cup W_2)}$ (all constructed via def \ref{def:G_marginalization} as just verified) coincide as CDMGs in the sense of def \ref{def-cdmg}, i.e.\ as four-tuples $(J, V, E, L)$. Concretely, the LN's displayed triple equality
    \[
        \lp G^{\sm W_1} \rp^{\sm W_2} \;=\; \lp G^{\sm W_2} \rp^{\sm W_1} \;=\; G^{\sm (W_1 \cup W_2)}
    \]
    is the conjunction of the two following equalities of CDMGs:
    \begin{enumerate}[label=(\alph*)]
        \item \emph{Composition with $W_1$ first, then $W_2$.}
              \[ \lp G^{\sm W_1} \rp^{\sm W_2} \;=\; G^{\sm (W_1 \cup W_2)}. \]
        \item \emph{Composition with $W_2$ first, then $W_1$.}
              \[ \lp G^{\sm W_2} \rp^{\sm W_1} \;=\; G^{\sm (W_1 \cup W_2)}. \]
    \end{enumerate}
    Together (a) and (b) yield the LN's ``symmetry under swap'' reading $\lp G^{\sm W_1} \rp^{\sm W_2} = \lp G^{\sm W_2} \rp^{\sm W_1}$ by transitivity of equality. The equality of CDMGs in (a) and (b) is understood componentwise on the four components $(J, V, E, L)$ of def \ref{def-cdmg}: the $J$- and $V$-component equalities are the carrier identifications already displayed above, while the $E$- and $L$-component equalities -- the substantive content of the lemma -- assert that the directed-walk-through-$W$ predicate $\Phi_E$ (item~iii.\ of def \ref{def:G_marginalization}) and the bifurcation-through-$W$ predicate $\Phi_L$ (item~iv.\ of def \ref{def:G_marginalization}), instantiated iteratively first with $W = W_1$ on $G$ and then with $W = W_2$ on $G^{\sm W_1}$ (respectively first with $W = W_2$ on $G$ and then with $W = W_1$ on $G^{\sm W_2}$), pick out exactly the same directed edges (over $\lp J \cup (V \sm (W_1 \cup W_2)) \rp \times (V \sm (W_1 \cup W_2))$) and the same bidirected edges (over $(V \sm (W_1 \cup W_2)) \times (V \sm (W_1 \cup W_2))$ with the $\ul{v} \neq \ol{v}$ constraint on $L$ from item~iv.) as the single instantiation of $\Phi_E$ and $\Phi_L$ with $W = W_1 \cup W_2$ on $G$. The conjunctive unpacking of (a) and (b) into the four field-by-field equalities of $J$, $V$, $E$, and $L$ is deferred to the proof.

    %This shows that the notation:
    %\[ G(V \sm (W_1 \cup W_2)|\doit(J))  \]
    %is unambiguous.
\end{Lem}

\end{document}
